\chapter{Elbow Pain}

\section*{Definition}
Pain or discomfort localized to the elbow joint and surrounding soft tissues, commonly arising from tendinopathy, arthritis, trauma, or nerve entrapment.

\section*{What Happens in Your Body}
Elbow pain localizes based on anatomy: lateral epicondylitis (tennis elbow) involves degenerative changes in the extensor carpi radialis brevis origin; medial epicondylitis (golfer's elbow) affects the common flexor origin; olecranon bursitis involves inflammation of the bursa overlying the olecranon process. Articular pain suggests osteoarthritis, inflammatory arthritis, or intra-articular loose bodies.

\section*{Causes}
\begin{itemize}
  \item Lateral epicondylitis (tennis elbow)
  \item Medial epicondylitis (golfer's elbow)
  \item Olecranon bursitis
  \item Elbow osteoarthritis
  \item Rheumatoid arthritis
  \item Cubital tunnel syndrome (ulnar nerve entrapment)
  \item Distal biceps or triceps tendinopathy
  \item Elbow dislocation or fracture (radial head, olecranon)
  \item Ligamentous injury (UCL tear — Tommy John lesion)
  \item Septic arthritis or gout
\end{itemize}

\section*{When to See a Doctor}
See your doctor if:
\begin{itemize}
  \item Symptoms are severe or getting worse over time
  \item Elbow pain after acute injury with deformity or inability to move, fever, or severe pain with swelling and warmth
  \item You have other concerning symptoms such as fever, weight loss, or bleeding
\end{itemize}

\section*{Related Symptoms}
\begin{itemize}
  \item Arm pain
  \item Forearm pain
  \item Numbness in fingers
  \item Swelling
  \item Decreased range of motion
\end{itemize}

\textbf{Important:} If this symptom persists, your doctor will check for serious underlying causes.

\section*{Self-Care / Home Management}
\begin{itemize}
  \item Rest the elbow and avoid repetitive motions that aggravate the pain (e.g., gripping, lifting).
  \item Apply ice packs for 15-20 minutes several times daily during the acute phase.
  \item Use a compression bandage or elbow brace to support the joint.
  \item Over-the-counter NSAIDs (ibuprofen, naproxen) can help reduce pain and inflammation.
  \item Gradually resume activity with proper technique and ergonomic adjustments.
\end{itemize}

\section*{How Your Doctor Will Evaluate You}
\begin{itemize}
  \item Your doctor will palpate specific anatomical landmarks (lateral and medial epicondyles, olecranon) to localize tenderness.
  \item Range of motion, strength, and neurovascular status of the elbow and forearm will be assessed.
  \item X-rays are obtained to evaluate for fractures, arthritis, or loose bodies.
  \item Ultrasound or MRI may be ordered to assess tendons, ligaments, and bursae.
  \item Referral to an orthopedic specialist or physical therapist for specialized management.
\end{itemize}

\section*{Treatment Options}
\begin{itemize}
  \item Physical therapy focusing on eccentric strengthening for lateral and medial epicondylitis.
  \item Corticosteroid injections for acute tendinopathy or olecranon bursitis.
  \item NSAIDs for pain and inflammation in arthritis or tendinopathy.
  \item Surgical intervention (release, debridement, or repair) for refractory cases.
  \item Activity modification and ergonomic assessment to prevent recurrence.
\end{itemize}


\section*{References}
\begin{thebibliography}{99}
\bibitem{elbow_pain:nirschl2011}
Nirschl RP, Ashman ES. Elbow tendinopathy: tennis elbow. \textit{Clin Sports Med}. 2003;22(4):813--836.
\bibitem{elbow_pain:shiri2015}
Shiri R, Viikari-Juntura E, Varonen H, Heli\"{o}vaara M. Prevalence and determinants of lateral and medial epicondylitis. \textit{Am J Epidemiol}. 2006;164(11):1065--1074.
\bibitem{elbow_pain:bisset2005}
Bisset L, Paungmali A, Vicenzino B, Beller E. A systematic review and meta-analysis of clinical trials on physical interventions for lateral epicondylalgia. \textit{Br J Sports Med}. 2005;39(7):411--422.
\bibitem{elbow_pain:calmbach2007}
Calmbach WL, Hutchens M. Evaluation of patients presenting with knee pain: part II. Differential diagnosis. \textit{Am Fam Physician}. 2003;68(5):917--922. [adapted for elbow]
\bibitem{elbow_pain:buchbinder2011}
Buchbinder R, Green SE, Struijs P. Tennis elbow. \textit{BMJ Clin Evid}. 2008;2008:1117.
\bibitem{elbow_pain:calfee2008}
Calfee RP, Manske PR, Gelberman RH, et al. Clinical assessment and treatment of elbow disorders. \textit{J Hand Surg Am}. 2008;33(8):1376--1384.
\end{thebibliography}
